\scanpage{6}
\mspage{4}

% pg 41
\tikz[baseline=(c.base)]{\node[draw,ellipse,inner sep=1pt] (c) {\uline{pg 41}};}\quad $l-2$\qquad ``fonction de Dirac'' est malheureux (sur un groupe, on ne sait trop si c'est \quad ou $\left\{\begin{array}{l}f(e)=1,\;0\ \text{ailleurs}\\f(e)=\#G,\;0\ \text{ailleurs}\end{array}\right.$ )

% left of pg 41
% slanted
\margin{guillemets suffisent ?}

% pg 42
\uline{pg 42}\quad début dém. 4.7 : le \unsure{cas|aps} n'est pas clair pour $n=2$, et je n'ai pas confiance en dehors du cas strictement local.

% division
\hrule\smallskip\hrule\smallskip\hrule

% section 5
\tikz[baseline=(c.base)]{\node[draw,ellipse,inner sep=1pt] (c) {5};}

% pg 45
\uline{pg 4\deleted{\unsure{7}}5} $l\,3$ ss\quad français : beaucoup de ``donc''.

% pg 46
\uline{pg 46}\quad Quel est le sens de (5.2.3,1) et de \deleted{\unsure{le|la}} ``par définition'' qui suit -- S'agit-il de pro-objets, ou de schémas.

% left of pg 46
\margin{\tikz[baseline=(c.base)]{\node[draw,ellipse,inner sep=1pt] (c) {je le dis !};}}

% pg 48
\uline{pg 48}\quad $l-4$\quad Pourquoi as-tu barré connexe ?

% pg 53
\deleted{\uline{pg 53}\quad 5.7 est une redite pour $l\ne p$, et un peu idiot pour $l=p$.}\quad (utilisé en 5.15)

% pg 56 bis
\tikz[baseline=(c.base)]{\node[draw,ellipse,inner sep=1pt] (c) {?\quad\uline{pg 56 bis}\quad 5.15.};}\quad tu n'a plus besoin ici de $\operatorname{car}(K)\ne l$

% 54 quater
54 quater -- 5.17.\qquad seraient plutôt que seront.

% 55 ter
\uline{55 ter}\quad $l-10$. Il serait honnête de dire que 6.6 utilise $\operatorname{car}(K)=0$, via Tate

% left of 55 ter
% slanted
\margin{\tikz[baseline=(c.base)]{\node[draw,ellipse,inner sep=1pt] (c) {\shortstack{c'est dit\\dans les \unsure{hyp.|hyp}\\de 5.13}};}}

% division
\hrule\smallskip\hrule\smallskip\hrule

% section 6, pg 57/58
\tikz[baseline=(c.base)]{\node[draw,ellipse,inner sep=1pt] (c) {6};}\quad \uline{pg 57/58} : tu supposes tantôt $A^0$ $l$-divisible (58, $l-11$), tantôt non (57, $l\,3$)

% pg 58
\uline{pg 58}\quad $l-5$\quad $({}_{p^\infty}X)^{\text{\tikz[baseline=(c.base)]{\node[draw,ellipse,inner sep=1pt] (c) {0};}}}$\quad \tikz[baseline=(c.base)]{\node[draw,ellipse,inner sep=1pt] (c) {?};}\qquad itou plus loin --

% pg 59
\uline{pg 59}\qquad Je ne comprends rien -- Que fais-tu de la partie discrète ? --

% left of pg 59
% slanted
\margin{\uline{CANULAR ?!}}

% pg 59, continuation
Je ne crois pas que la plus grande \deleted{\unsure{s|ss|de}} sous VA ordinaire puisse se relever (\~m en cas de bonne réduction).

% last lines
\uline{i e}\quad 63 est faux\quad (déformer de façon générale $E\times F$, $E$ et $F$ elliptiques, $E$ ordinaire, $F$ non ordinaire ; général $\Rightarrow$ générique simple.
% archive sheet number: 5
